\documentclass[12pt,a4paper]{article}

\usepackage[T2A]{fontenc}
\usepackage[utf8]{inputenc}
\usepackage[russian]{babel}
\usepackage{amsmath,amssymb}
\usepackage{graphicx}
\usepackage{geometry}
\geometry{margin=2.5cm}

\begin{document}

\vspace{2cm}

\section*{\centering Скобка чётности по движениям на графах}

\vspace{2cm}

\section*{1. Определения}

\subsection*{Определение 1: Разведение в множестве вершин}

Пусть $G$ --- граф, $S \subseteq V(G)$, $|S| = k$, и вершины пронумерованы так, что сначала идут вершины из $S$, затем вершины из $V(G)\setminus S$.

Говорим, что граф $H$ получен разведением в множестве вершин $S$, если существуют коэффициенты $\alpha=(\alpha_1,\ldots,\alpha_k)\in\{0,1\}^k$ такие, что $H$ изоморфен графу с матрицей смежности $A'$, полученной следующим образом:
\begin{enumerate}
\item Формируем матрицу
\[
M = A(G) + \operatorname{diag}(\alpha_1,\ldots,\alpha_k,0,\ldots,0).
\]
\item К матрице $M$ разрешено применять только элементарные преобразования третьего типа над $\mathbb{F}_2$: прибавление $i$-й строки к $j$-й или прибавление $i$-го столбца к $j$-му.
\item Разбиение $M$ на блоки задаётся как
\[
M=
\begin{pmatrix}
D & B\\
B^\top & C
\end{pmatrix},
\]
где $D$ --- левый верхний блок размера $k\times k$. Далее выполняем строго следующую процедуру:
\begin{enumerate}
\item сначала, используя только первые $k$ строк и первые $k$ столбцов, приводим блок $D$ к $I_k$;
\item затем с помощью преобразований строк зануляем левый нижний блок (при этом правый нижний блок может изменяться);
\item в конце преобразованиями столбцов зануляем правый верхний блок.
\end{enumerate}
В результате матрица приводится к виду
\[
\begin{pmatrix}
I_k & 0\\
0 & A'
\end{pmatrix},
\]
где $A'$ --- симметричная матрица.

\textbf{Замечание о диагонали.}
Элементы главной диагонали матрицы $A'$ (дополнения Шура) не рассматриваются:
граф $H_{S,\alpha}$ строится только по внедиагональным элементам матрицы $A'$.
Иными словами, ребро $\{i,j\}$ добавляется тогда и только тогда, когда $i\neq j$ и $A'_{ij}=1$.
\end{enumerate}

\textbf{Критерий возможности разведения:}
вводим подматрицу смежности $A(S)$ размера $k\times k$, индуцированную на $S$, и проверяем условие
\[
\det\bigl(A(S)+\operatorname{diag}(\alpha_1,\ldots,\alpha_k)\bigr)\neq 0 \pmod{2}.
\]

\subsection*{Лемма (однозначность процедуры)}

Пусть для фиксированных $G$, $S$ и $\alpha$ выполнено
\[
\det\bigl(A(S)+\operatorname{diag}(\alpha_1,\ldots,\alpha_k)\bigr)\neq 0 \pmod{2}.
\]
Тогда описанная выше процедура приводит матрицу $M$ к блочному виду
\[
\begin{pmatrix}
I_k & 0\\
0 & A'
\end{pmatrix}
\]
однозначно: итоговый блок $A'$ не зависит от выбора промежуточных преобразований.

\textit{Идея доказательства.}
В блочной записи $M=\bigl(\begin{smallmatrix}D & B\\ B^\top & C\end{smallmatrix}\bigr)$ после первого шага имеем $D=I_k$.
Дальнейшие шаги эквивалентны вычислению дополнения Шура по блоку $D$, поэтому
\[
A' = C + B^\top D^{-1}B
\quad \text{(над }\mathbb{F}_2\text{)}.
\]
Поскольку правая часть определяется только исходной матрицей $M$ (при фиксированных $S$ и $\alpha$), итоговый блок $A'$ однозначен.

\subsection*{Определение 2: Скобка чётности}

Вершина $v\in V(G)$ называется \emph{чётной} (в обычном смысле), если её степень $\deg v$ чётна.
Обозначим $V_{\mathrm{even}}(G)=\{v\in V(G)\mid \deg v\equiv 0\pmod{2}\}$.

Скобкой чётности графа $G$ называется формальная линейная комбинация над полем $\mathbb{F}_2$:
\[
[G]=\sum_{(S,\alpha)\ \text{допустимо}} [H_{S,\alpha}]_{\mathrm{II}},
\]
где сумма берётся по всем допустимым разведениям $(S,\alpha)$ только для $S=V_{\mathrm{even}}(G)$, а $[H_{S,\alpha}]_{\mathrm{II}}$ --- класс эквивалентности графа $H_{S,\alpha}$ относительно движений Рейдемейстера второго типа.

\end{document}
